\section{Cross-Domain Agentic Visual Generation}
\label{sec:cross-domain}

This chapter examines agentic visual generation across domains whose artifacts, interfaces, and evaluation signals differ. The comparison follows the operating loop introduced earlier: the retained state identifies what the system can address, the selected action changes an artifact or execution process, and the available evidence determines the next operation. The cases cover multimodal storytelling, long-horizon multimodal interaction, game-development evaluation, node-based editing, embodied navigation videos, personalized training media, and video-generation infrastructure. They show how the same control questions are instantiated when the system coordinates several media, exposes an editable graph, uses a visual trajectory to guide physical action, checks human constraints, or optimizes deployment.

\paragraph{Multimodal Qinqiang opera generation.}
A multi-agent Qinqiang opera framework converts a source story into a culturally grounded script, stage-scene descriptions and images, and synthesized vocal performance. One agent generates acts, scenes, dialogue, stage directions, and performance cues; a second converts script descriptions into visual scenarios; and a third assigns speech-synthesis parameters to dialogue and lyrics. The case study on \emph{Dou E Yuan} reports higher expert ratings than a single-agent baseline, while ablations associate the visual agent with visual coherence and the speech agent with speech accuracy~\cite{avg250415552}.

\paragraph{Cognitive-structured multimodal interaction.}
The Cognitive-structured Multimodal Agent addresses long-horizon sessions that combine visual understanding, image generation, editing, composition, topic switching, and references to earlier images. A Perceptual Abstraction Engine stores image tags, descriptions, and thumbnails in episodic visual memory; a Cognitive Retrieval Engine selects relevant episodes; and a Multimodal Executive Controller dispatches the retrieved context to generation, editing, composition, or response modules. The reported experiments show improved visual retrieval and generation quality in 20-turn sessions through selective visual-context recall~\cite{avg260708497}.

\paragraph{Game-development evaluation.}
GameDevBench evaluates coding agents on Godot tasks that combine repository-scale code with sprites, shaders, animations, scene graphs, UI layouts, physics, cameras, and gameplay behavior. Its tasks provide reference solutions and deterministic tests, while optional editor screenshots and runtime videos expose visual and temporal state to evaluated coding agents. The benchmark shows that multimodal observation improves pass@1 on some configurations, while recurring failures involve node hierarchy, materials, layout, animation state, camera framing, and asset selection~\cite{avg260211103}.

\paragraph{Node-Based Editing.}
Node-Based Editing represents text, audio, images, and video as editable nodes in a multimodal graph. A task-selection agent routes requests among story generation, node-structure reasoning, diagram formatting, context generation, and node editing. Users can edit a node, branch a story, or regenerate a selected node while retaining the rest of the graph. The node graph therefore provides an intermediate representation for passing context between modalities and for limiting regeneration to the affected part of a multimodal artifact~\cite{avg251103227}.

\paragraph{Action Agent.}
Action Agent separates language-guided navigation into trajectory imagination and low-level execution. In the first stage, an orchestration agent selects a video-generation pipeline, revises prompts through structured validation, stores successful strategies across tasks, and accepts a first-person navigation video when prompt adherence, physical plausibility, and visual quality meet the reported criteria. FlowDiT then maps the validated video and language instruction to continuous velocity commands for different robot embodiments. The paper reports 86\% success for generated navigation videos across 50 tasks, 73.2\% navigation success in simulation, and 64.7\% task completion on a real Unitree G1 under open-loop execution~\cite{avg260501477}.

\paragraph{Agentic AI for Personalized Physiotherapy.}
Agentic AI for Personalized Physiotherapy organizes a rehabilitation pipeline around four cooperating agents. The Clinical Extraction Agent converts medical notes into kinematic constraints; the Video Synthesis Agent uses those constraints to create a patient-specific exercise demonstration; the Vision Processing Agent estimates body pose from camera frames; and the Diagnostic Feedback Agent compares the estimated motion with the constraints and issues corrective instructions. A shared patient state carries the prescription, extracted limits, pose, and feedback through the loop. The paper presents a prototype architecture, implementation description, and preliminary component estimates, so its contribution is a design feasibility result for a safety-oriented generative-video workflow~\cite{avg260421154}.

\paragraph{Sol Video Inference Engine.}
Sol Video Inference Engine applies agentic coordination to the inference stack surrounding video generation. For a specified model, hardware platform, and serving configuration, parallel skill agents optimize caching, sparse attention, token pruning, quantization, and kernel fusion. An integrator composes the candidate techniques into a full stack, and a human validator checks whether the resulting speedup preserves visual quality. Across Cosmos3-Super, LTX-2.3, and SANA-Video, the framework reports more than 2$\times$ end-to-end acceleration with near-lossless VBench quality~\cite{avg260623743}.
